\section{Conormal layers and the contraction interfaces}\label{sec:contraction-inputs}

We next prove the ordinary conormal statement used to extend a contraction
of one prime floor to an ambient contraction.  It concerns divisorial
ideals and their actual reflexive powers.  We then construct analytic
base thickenings from the resulting direct images and apply the analytic
blowdown criterion to obtain the ambient contraction.

\begin{proposition}[Conormal layers of a prime floor]\label{prop:conormal-layers}
Let \((X,S+B)\) be an ordinary plt pair, where \(X\) is a normal connected
compact K\"ahler space, \(S\) is an irreducible effective \(\Q\)-Cartier
prime, \(B\geq0\) is a rational \(\Q\)-Cartier divisor not containing
\(S\), and \(K_X+S+B\) is an actual \(\Q\)-Cartier adjoint.  Suppose
\(\dim S\leq3\).  Write the actual residue adjunctions as
\[
 (K_X+S)|_S=K_S+\Theta,\qquad
 (K_X+S+B)|_S=K_S+B_S .
\]
Let \(\phi:S\to W\) be a projective surjection to a normal compact
analytic space with \(\phi_*\OO_S=\OO_W\).  Suppose that the actual
rational lines \(-(K_S+B_S)\) and \(-S|_S\) are \(\phi\)-ample.
For \(k\geq0\), put
\[
 \mathscr I_k=\OO_X(-kS),\qquad
 \mathscr E_k=\mathscr I_k/\mathscr I_{k+1}.
\]
Then \(S\) is normal, \((S,B_S)\) is klt, and each \(\mathscr E_k\)
is a rank-one reflexive sheaf on \(S\).  For any positive global Cartier
index \(q\) of \(S\), there is a finite effective rational Weil divisor
\(\Xi_k\) on \(S\), with \(q\Xi_k\) integral, such that
\begin{equation}\label{eq:conormal-actual-identity}
 \mathscr E_k^{[q]}
 \simeq \bigl(\OO_X(-kqS)|_S\bigr)\otimes\OO_S(-q\Xi_k),
 \qquad 0\leq\Xi_k\leq\Theta\leq B_S .
\end{equation}
Here \(\mathscr E_k^{[q]}=(\mathscr E_k^{\otimes q})^{**}\) is formed
on \(S\), and \(\OO_S(-q\Xi_k)\) is a divisorial reflexive sheaf,
which need not be invertible.  Moreover
\begin{equation}\label{eq:conormal-vanishing}
 R^i\phi_*\mathscr E_k=0\qquad(i>0,\ k\geq1).
\end{equation}
\end{proposition}
\begin{proof}
Choose an integer \(N>0\) clearing the actual adjoint index and the
Cartier indices of \(S+B\).  On the common codimension-one locally free
open, and hence everywhere by reflexive extension,
\[
 \omega_X^{[N]}
 \simeq \OO_X\bigl(N(K_X+S+B)\bigr)\otimes\OO_X\bigl(-N(S+B)\bigr).
\]
Thus this reflexive canonical power is an actual line.  In particular
\(K_X+S\) is an actual rational adjoint; no global canonical Weil
generator is being chosen.  Subtracting the effective \(\Q\)-Cartier
boundary from discrepancies shows that \((X,S)\) is plt and \(X\) is
klt.  Ordinary plt adjunction
\cite[Lemma~5.3]{DasHacon2024OnMMP} gives normal \(S\) and klt
\((S,B_S)\).  With compatible residue embeddings, restriction of the
canonical section of a Cartier multiple of \(B\) gives
\[
 B_S=\Theta+B|_S,\qquad B|_S\geq0.
\]
The coefficient computation below also proves \(\Theta\geq0\).
The notation \(S|_S\) always denotes the restriction of the actual
rational normal line.

\subsubsection*{The analytic quotient in codimension two}

Fix a prime \(P\) on \(S\).  Near a general point choose a holomorphic
equation \(v\) for a Cartier multiple \(mS\), and normalize the full
finite root cover
\[
 t^m=v.
\]
The root algebra is reduced before normalization.  It is free with basis
\(1,t,\ldots,t^{m-1}\) over the normal local domain.  Its generic
Kummer algebra is reduced in characteristic zero because \(v\ne0\),
and freeness makes the map to that generic algebra injective.  A
nilpotent element must therefore be zero.
Finite analytic normalization
\cite[Part~B, Section~4, Corollaries~2--3]{Houzel1961}
gives a finite normal cover \(\tau:Y\to U\) with its complete
\(\mu_m\)-action.  It can be disconnected, and we retain all components.
The invariant subalgebra of its normalization is \(\OO_U\):
in the total meromorphic algebra the invariants are the meromorphic
field of \(U\), and normality of \(U\) identifies its integral elements.
Thus the quotient is the original normal space.

This cover is étale at every codimension-one point.  Off \(S\), the
equation is a root of a unit.  At a general smooth point of \(S\), write
\(v=u x^m\) for a unit \(u\); after choosing a holomorphic root of \(u\),
normalization is a disjoint union of the branches
\(t=\zeta u^{1/m}x\).  The lifted floor
\[
 S_Y=\operatorname{div}(t)
\]
is reduced and Cartier.  Its prime coefficients are one, and on a
normal space the principal divisorial ideal with these coefficients is
the radical ideal of their union.

The finite discrepancy formula is valid here in the analytic category.
Normalize a base change of a model carrying a tested divisorial valuation.
The Hurwitz formula for the corresponding discrete valuation rings and
the actual identity \(K_Y+S_Y=\tau^*(K_X+S)\) give
\begin{equation}\label{eq:cover-discrepancy}
 a(F;Y,S_Y)=e(F/E)\,a(E;X,S).
\end{equation}
Extensions of the valuation of \(S\) are exactly the height-one
valuations of the lifted floor.  All exceptional log discrepancies
remain positive, so the lifted pair is plt.  Applying the same formula
without the floor shows that the components of \(Y\) are klt.

The lifted floor is normal.  One can use the ordinary plt adjunction just
cited, or the following independent connectedness argument.  On a
projective SNC resolution of the lifted pair, its coefficient-one locus
is the union of the strict floor components, with no exceptional
component.  Two such components cannot meet: blowing up their
intersection would give an exceptional log discrepancy zero.  They
are therefore disjoint and smooth.  Lemma~\ref{lem:floor-connectedness}
gives the equality of their direct-image structure sheaf with that of
the reduced floor.  Factoring the map through the finite normalization
of the floor, and applying normal bimeromorphic Hartogs extension on
that normalization, identifies this direct image with the normalization
sheaf.  The equality forces the normalization to be an isomorphism.
This use of the connectedness lemma is independent of any contraction
existence assertion.

Choose a component \(P'\) above \(P\) and a general point on it.
The normal \(S_Y\) is smooth there, since \(P'\) has codimension one
in it, and \(P'\) is smooth after a further generic choice.  The ambient
\(Y\) is smooth there as well.  Indeed, in its local ring \(R'\) the
nonzerodivisor \(t\) has regular quotient \(R'/(t)\).  Lifting its
minimal generators and adjoining \(t\) bounds the embedding dimension
of \(R'\) by \(\dim R'\), so \(R'\) is regular.

Remove the finitely many proper fixed loci on \(P'\).  At a remaining
point its stabilizer \(G\subset\mu_m\) fixes \(P'\) pointwise.  Holomorphic
linearization can preserve the smooth flag
\(P'\subset S_Y\subset Y\): average lifts of an eigenbasis of the
cotangent space inside the respective invariant ideals, and use their
independent differentials as local coordinates.  The coordinates tangent
to \(P'\) are fixed.  The stabilizer has no ineffective element: its
action sends \(t\) to \(\zeta t\), and the nonzerodivisor \(t\) forces
\(\zeta=1\) for an element acting trivially on the local germ.  A
nonidentity stabilizer element cannot fix a transverse hyperplane,
because \(\tau\) is étale in codimension one.  Both transverse
characters are consequently faithful.  If \(r=|G|\), rescaling the
generator gives the analytic germ
\begin{equation}\label{eq:plt-analytic-quotient}
 (X,S,P)\simeq
 \bigl((\C^2,\{x=0\},0)/\mu_r(1,a)\bigr)
      \times\C^{\dim X-2},\qquad \gcd(a,r)=1 .
\end{equation}
The smooth case is \(r=1\).  The local Cartier index of \(S\) is \(r\):
\(x^r\) descends, whereas invariance of a unit times \(x^d\), evaluated
at the fixed point, forces \(r\mid d\).

On the quotient floor let \(z=y^r\).  Residue of the \(r\)-th
log-canonical power changes \((dy)^{\otimes r}\) into a nonzero constant
times \(z^{-(r-1)}(dz)^{\otimes r}\); the tangential factors are
unchanged.  The actual meromorphic residue embedding therefore gives
\begin{equation}\label{eq:bare-different-index}
 \operatorname{coeff}_P\Theta=\frac{r-1}{r}.
\end{equation}
This derivation is on the analytic cover and does not infer a
higher-dimensional analytic quotient theorem from an algebraic slice.

\subsubsection*{Depth and the canonical multiplication}

The sheaves \(\mathscr I_k\) are maximal Cohen--Macaulay on \(X\).
Here is the index-cover argument at every point of \(X\).  On a small
open about an arbitrary point, repeat the full normalized root construction
using a local equation for a Cartier multiple \(mS\).  Codimension-one
étaleness and the discrepancy calculation above do not require the
generic choice on \(P\).  The klt components of this full cover \(Y\)
are Cohen--Macaulay by the last assertion of
Lemma~\ref{lem:floor-connectedness}; its proof uses relative vanishing,
canonical torsion-freeness, and coherent duality.  The total meromorphic
Kummer algebra has basis \(1,t,\ldots,t^{m-1}\), even when it splits.
Its character-\(j\) holomorphic summand on this open is
\[
 \OO_X(jS)t^j\qquad(0\leq j<m).
\]
Indeed regularity in each height-one discrete valuation ring requires
the coefficient to have order at least \(-j\) along \(S\) and at least
zero elsewhere.  Normality supplies the equality of these sheaves.
Up to tensoring by the actual Cartier line of a multiple of \(mS\),
these are all the \(\mathscr I_k\).  Averaging over \(\mu_m\) splits
them as modules.

For the depth assertion, let \(R=\OO_{X,x}\) have dimension \(d\) and
choose a system of parameters.  Every normalization component dominates
the base open: the generic Kummer algebra is a separable product of
fields, and the full normalization retains all of those components.
Thus each local ring of \(Y\) above \(x\) has dimension \(d\), and
finiteness makes the extended parameter ideal primary for its maximal
ideal.  Cohen--Macaulayness upstairs
makes the same parameter sequence regular.  Testing at all upstairs
maximal ideals makes it regular on the finite semilocal \(R\)-module
\((\tau_*\OO_Y)_x\).  It remains regular on each split summand, and
hence on \(\mathscr I_k\).  No flatness of the finite cover is used.

Divisor orders give
\(\mathscr I_j\mathscr I_\ell\subset\mathscr I_{j+\ell}\); also
\(\mathscr I_1\) is the reduced ideal of \(S\).  Thus
\(\mathscr E_k\) is a coherent sheaf on \(S\).  The depth lemma in
\[
 0\longrightarrow\mathscr I_{k+1}\longrightarrow\mathscr I_k
 \longrightarrow\mathscr E_k\longrightarrow0
\]
gives depth at least \(\dim X-1\) along \(S\).  The quotient has
generic rank one on the irreducible \(S\), so its support is all of
\(S\) and has local dimension \(\dim X-1\).  Depth cannot exceed
this dimension, and equality follows.  Thus it is Cohen--Macaulay on
\(S\), has no embedded associated
component, and is torsion-free and \(S_2\).  Normality of \(S\) makes
it rank-one reflexive.

In \eqref{eq:plt-analytic-quotient}, exactness of finite-group invariants
identifies
\[
 \mathscr E_k=(x^k\C\{y,\mathbf u\})^{\mu_r},
\]
where \(\mathbf u\) denotes the fixed tangential coordinates.  Let
\(j_{k,P}\) be the unique integer with
\[
 0\leq j_{k,P}<r,\qquad k+a j_{k,P}\equiv0\pmod r .
\]
Each invariant series factors as
\(x^k y^{j_{k,P}}h(y^r,\mathbf u)\).  Hence this is a free rank-one
module at the general point of \(P\), and its \(r\)-th multiplication
into the degree \(kr\) layer is
\[
 (x^k y^{j_{k,P}})^r=(x^r)^k z^{j_{k,P}} .
\]
Relative to the local frame \((x^r)^k\) of the actual line
\(\OO_X(-krS)|_S\), its vanishing order is \(j_{k,P}\).  By
\eqref{eq:bare-different-index},
\begin{equation}\label{eq:conormal-local-defect}
 \xi_{k,P}:=\frac{j_{k,P}}r
 \quad\hbox{satisfies}\quad
 0\leq\xi_{k,P}\leq\operatorname{coeff}_P\Theta .
\end{equation}

Now choose the global Cartier index \(q\) in the statement.  Multiplication
of the divisorial filtration defines a canonical global map
\begin{equation}\label{eq:conormal-multiplication}
 \mu_k:\mathscr E_k^{\otimes q}\longrightarrow\mathscr E_{kq}
     =\OO_X(-kqS)|_S=:\mathscr A_k .
\end{equation}
Replacing one factor by \(\mathscr I_{k+1}\) raises its product into
\(\mathscr I_{kq+1}\), so the map is well defined.  The equality on the
right follows from
\(\mathscr I_{kq+1}=\mathscr I_1\otimes\OO_X(-kqS)\).
The image \(\mathscr J_k\) is a coherent rank-one subsheaf of the actual
line \(\mathscr A_k\).  Its reflexive hull is that line tensored with
the divisorial ideal of an effective integral Weil divisor.  At every
codimension-one point the source of \(\mu_k\) is free and the map is
injective.  Thus its reflexive hull is canonically the same
\(\mathscr J_k^{**}\).

The local index \(r\) divides \(q\), so the order of \(\mu_k\) at \(P\)
is \(qj_{k,P}/r\).  Define
\(\Xi_k=\sum_P\xi_{k,P}P\).  It is finite, either from the coherent
image on the compact \(S\) or from the bound by \(\Theta\leq B_S\).
Taking the reflexive image in \eqref{eq:conormal-multiplication} gives
exactly \eqref{eq:conormal-actual-identity}.  The canonical map into
the specified target line supplies the global identity and excludes an
undetermined line-bundle factor.

\subsubsection*{Vanishing on a small floor model}

We use the projective relative MMP in dimension at most three to obtain
a small ordinary \(\Q\)-factorial model \(p:S'\to S\).  For precision,
take a projective log resolution of \((S,B_S)\), principalizing the
coherent ideal of its finite boundary support; no individual
\(\Q\)-Cartier assumption on \(B_S\) is required.  Give each exceptional
prime an effective rational coefficient below one and strictly above
its crepant coefficient.  For the resulting effective SNC klt boundary
\(C_{\widehat S}\),
\[
 K_{\widehat S}+C_{\widehat S}
   =\rho^*(K_S+B_S)+F,\qquad F\geq0
\]
with \(F\) exceptional and positive on every exceptional prime.  This
is the actual meromorphic adjoint identity.  The adjoint is relatively
pseudo-effective.  The morphism is a projective surjection between normal
compact analytic spaces, the source is smooth and ordinary
\(\Q\)-factorial, and its dimension is at most three.  Thus
\cite[Proposition~2.26]{DasHacon2024} applies.  On its relative minimal
model the transform of \(F\) is exceptional and relatively nef, so
negativity \cite[Lemma~1.3]{Wang2021} makes it zero.  No step extracts
a prime, hence the resulting \(p:S'\to S\) is small and projective.
The source is ordinary globally \(\Q\)-factorial and compact K\"ahler,
and codimension-one comparison gives the actual crepant identity
\begin{equation}\label{eq:conormal-small-crepant}
 K_{S'}+B'_S=p^*(K_S+B_S)
\end{equation}
with the strict boundary and a klt pair.  This uses the relative
projective MMP as an external input; it is separate from the global
contraction construction considered here.

The finite strict transform \(\Xi'_k\) is \(\Q\)-Cartier by ordinary
global \(\Q\)-factoriality, as is \(B'_S\).  Put
\(B'_k=B'_S-\Xi'_k\).  It is effective by
\eqref{eq:conormal-actual-identity}.  Pullback of the effective
\(\Q\)-Cartier \(\Xi'_k\) lowers the crepant coefficients of the klt
pair, so \((S',B'_k)\) is klt.

Work over a connected relatively compact Stein open \(U\subset W\) whose closure
lies in a chosen larger open.  Properness and connected fibers make
\(S_U\) connected, and normality then makes it irreducible.  Relative
generation for a large twist
by a \(\phi\)-ample line \cite[Theorem~2.12]{DasHaconPaun2024},
followed by Cartan generation on the Stein
base, gives a meromorphic section of \(\mathscr E_k|_{S_U}\): take the
quotient of a nonzero section of the twist and a nonzero section of the
twisting line.  The same argument applies to the actual line
\(\mathscr A_k\).  Write
\[
 \mathscr E_k|_{S_U}=\OO_{S_U}(G_k),\qquad
 \mathscr A_k|_{S_U}=\OO_{S_U}(H_k)
\]
for the resulting integral Weil divisor \(G_k\) and Cartier divisor
\(H_k\).  Applying the actual identity
\eqref{eq:conormal-actual-identity} to these meromorphic sections gives
a genuine principal relation
\[
 qG_k\sim H_k-q\Xi_k .
\]
Let \(G'_k\) be the strict transform on \(S'_U\).  A small map has no
local exceptional prime either.  Strict transform therefore takes
\(H_k\) to its Cartier pullback and a principal divisor to that of the
pulled-back meromorphic function.  It follows that
\begin{equation}\label{eq:conormal-local-cartier}
 qG'_k\sim p^*H_k-q\Xi'_k .
\end{equation}
After clearing the global indices of \(\Xi'_k\), the right side is
Cartier.  Hence the locally chosen integral \(G'_k\) is
\(\Q\)-Cartier.  This local assertion follows from the displayed
identity; ordinary global \(\Q\)-factoriality supplied the index of the
globally defined \(\Xi'_k\).
The actual adjoint in \eqref{eq:conormal-small-crepant}, together with
the globally \(\Q\)-Cartier boundary \(B'_S\), also makes a reflexive
canonical power on \(S'\) an actual line.  The same local generation
argument supplies a meromorphic canonical generator if a local Weil
representative is needed for the vanishing statement.

Equations \eqref{eq:conormal-small-crepant} and
\eqref{eq:conormal-local-cartier} give
\begin{equation}\label{eq:conormal-kv-difference}
 G'_k-(K_{S'}+B'_k)\sim_\Q p^*H,\qquad
 H=-(K_S+B_S)-kS|_S .
\end{equation}
The right side pulls back the defined full adjoint and normal rational
lines, and the locally chosen \(G'_k\) is \(\Q\)-Cartier by
\eqref{eq:conormal-local-cartier}.  The line \(H\) is \(\phi\)-ample.
Projective analytic
morphisms compose over a neighborhood of a compact subset of the final
base \cite[Remark~2.11]{DasHaconPaun2024}.  Here \(W\) is compact, so
\(\phi p\) is projective over the whole base; equivalently one may use
the compact closures in the local argument.  Both \(p\) and \(\phi p\)
are proper surjections.

The difference in \eqref{eq:conormal-kv-difference} is nef and big for
both maps in the convention of
\cite[Definition~2.40]{DasHaconPaun2024}.  Nefness follows by projecting
vertical curves: a curve for \(\phi p\) maps to a \(\phi\)-vertical
curve or a point, and the degree for \(p\) is zero.  For the relative
Iitaka condition, clear the index and write \(\mathscr H=\OO_S(hH)\)
for a \(\phi\)-ample line.  Normality and proper bimeromorphicity give
\(p_*\OO_{S'}=\OO_S\), so projection formula gives
\[
 (\phi p)_*p^*\mathscr H^m=\phi_*\mathscr H^m,\qquad
 p_*p^*\mathscr H^m=\mathscr H^m .
\]
For large \(m\), the first relative complete system is \(p\) followed
by the relative embedding from \(\mathscr H^m\).  Its image has relative
dimension \(\dim S'-\dim W\).  The second has image the base \(S\), of
relative dimension zero, equal to \(\dim S'-\dim S\).  These are exactly
the two maximal relative Iitaka dimensions.  For \(p\), relative bigness
is this zero-relative-dimension condition; the pulled line has degree zero
on its exceptional curves.

Analytic Kawamata--Viehweg vanishing
\cite[Theorem~2.41]{DasHaconPaun2024} now applies to the effective
rational klt boundary \(B'_k\), the integral \(\Q\)-Cartier divisor
\(G'_k\), and each of the two proper surjections.  It yields, locally
over \(U\),
\[
 R^ip_*\OO_{S'}(G'_k)=0,\qquad
 R^i(\phi p)_*\OO_{S'}(G'_k)=0\quad(i>0).
\]
The bijection of prime valuations for the small map, together with
normal divisorial extension, identifies
\[
 p_*\OO_{S'}(G'_k)=\OO_S(G_k)=\mathscr E_k .
\]
This is an equality of meromorphic sections satisfying the same
codimension-one order conditions.  Leray gives
\eqref{eq:conormal-vanishing} on \(U\), hence everywhere.  This proves
the proposition.
\end{proof}

For \(k\geq1\), put \(S_k=V(\mathscr I_k)\).  The inclusions
\(\mathscr I_1^k\subset\mathscr I_k\subset\mathscr I_1\) show that
\(\sqrt{\mathscr I_k}=\mathscr I_1\), so these thickenings have the
common underlying floor.  Their conormal sequence is
\[
 0\longrightarrow\mathscr E_k\longrightarrow\OO_{S_{k+1}}
 \longrightarrow\OO_{S_k}\longrightarrow0 .
\]
Use the underlying continuous map \(|\phi|:|S|\to|W|\) to put
\(\mathscr B_k=|\phi|_*\OO_{S_k}\).  Then
\(\mathscr B_1=\OO_W\), and \eqref{eq:conormal-vanishing} gives
epimorphisms of sheaves of rings on \(|W|\)
\begin{equation}\label{eq:thickening-ring-surjection}
 \mathscr B_{k+1}\twoheadrightarrow\mathscr B_k\qquad(k\geq1).
\end{equation}
This is the immediate cohomological assertion used in
\cite[Claim~5.9]{DasHacon2024OnMMP}.  Exactness is stalkwise, and no
splitting of the augmentation \(\mathscr B_k\to\OO_W\) is assumed.
The next construction realizes these ringed spaces analytically and
constructs the corresponding maps from \(S_k\).

\subsection{Analytic realization of the floor thickenings}
\label{subsec:analytic-thickenings}

Retain the divisorial ideals \(\mathscr I_k\), the thickenings \(S_k\),
the map \(\phi:S\to W\) from Proposition~\ref{prop:conormal-layers},
and the sheaves \(\mathscr B_k\) in
\eqref{eq:thickening-ring-surjection}.  We construct their analytic
structure from the already analytic source thickenings.

\begin{lemma}[Analytic base thickenings]\label{lem:analytic-thickenings}
For every \(k\geq1\), put
\[
 W_k=(|W|,\mathscr B_k).
\]
Then \(W_k\) is a complex space with canonical reduction \(W\), and
the canonical sheaf evaluation defines a proper surjective holomorphic
map \(\phi_k:S_k\to W_k\) with
\[
 (\phi_k)_*\OO_{S_k}=\OO_{W_k}.
\]
Its reduction is \(\phi\).  The quotients
\(\mathscr B_{k+1}\twoheadrightarrow\mathscr B_k\) define closed
analytic embeddings \(W_k\hookrightarrow W_{k+1}\).  They commute
with the embeddings \(S_k\hookrightarrow S_{k+1}\) and the maps
\(\phi_k\).
\end{lemma}
\begin{proof}
\par\medskip\noindent\textit{Canonical rings and their successive kernels.}\par
All direct images defining \(\mathscr B_k\) first use the continuous
map \(|\phi|:|S|\to|W|\).  The equality
\(\mathscr B_1=\OO_W\) is the canonical equality for the holomorphic
map \(\phi\).  Apply derived direct image for abelian sheaves to the
conormal sequence.  Its derived image on an \(\OO_S\)-module is the
underlying sheaf of the analytic higher direct image.  To see this
formally, an injective \(\OO_S\)-module is flasque: for open sets
\(V\subset U\), the extension-by-zero modules satisfy
\(j_{V!}\OO_V\hookrightarrow j_{U!}\OO_U\).  Adjunction identifies
\(\operatorname{Hom}_{\OO_S}(j_{U!}\OO_U,J)\) with
\(\Gamma(U,J)\), so injectivity of \(J\) makes the restriction of
sections surjective.  Flasque abelian sheaves are acyclic for continuous
direct image.  An injective
module resolution therefore computes both the analytic module higher
images and their underlying abelian-sheaf higher images.  Thus
\eqref{eq:conormal-vanishing} gives exact sequences
\begin{equation}\label{eq:analytic-thickening-sequence}
 0\longrightarrow \mathscr F_k:=\phi_*\mathscr E_k
 \longrightarrow\mathscr B_{k+1}
 \longrightarrow\mathscr B_k\longrightarrow0\qquad(k\geq1).
\end{equation}
The displayed quotients are ring maps; exactness here concerns their
underlying abelian sheaves.  Proper direct-image coherence for the
already holomorphic \(\phi\)
\cite[Theorem~1.1, p.~15-01]{Grothendieck1961TechniquesVIII}
makes \(\mathscr F_k\) a coherent
\(\OO_W\)-module.

In \(\OO_{S_{k+1}}\), the ideal \(\mathscr E_k\) is square-zero and
is annihilated by \(\mathscr I_1\): use
\(\mathscr I_k^2\subset\mathscr I_{2k}\subset\mathscr I_{k+1}\)
and \(\mathscr I_1\mathscr I_k\subset\mathscr I_{k+1}\).
Consequently the action of \(\mathscr B_{k+1}\) on the ideal
\(\mathscr F_k\) factors through \(\mathscr B_1=\OO_W\).
This is the \(\OO_W\)-module structure on the square-zero ideal used
in the following construction.

The composites of \eqref{eq:analytic-thickening-sequence} are
epimorphisms \(\mathscr B_k\to\OO_W\).  Their kernels are
\[
 \mathscr N_k=|\phi|_*(\mathscr I_1/\mathscr I_k),
 \qquad \mathscr N_k^k=0,
\]
by left exactness and \(\mathscr I_1^k\subset\mathscr I_k\).
Since \(W\) is reduced, \(\mathscr N_k\) is exactly the nilradical.
Each stalk of \(\mathscr B_k\) is local: an element lifting a unit in
\(\OO_{W,w}\) has an inverse by a finite geometric series in the
nilpotent kernel, whereas an element mapping to the maximal ideal is
not a unit.  Its residue field is \(\C\).  Thus \(W_k\) is already
a locally ringed space, and it remains to construct local analytic
presentations for it.

\par\medskip\noindent\textit{A local analytic presentation from source functions.}\par
Fix \(k\) and \(w\in W\).  Choose a local closed analytic embedding
of a neighborhood of \(w\) into an open polydisc.  Lift the finitely
many germs of its coordinate restrictions through
\(\mathscr B_k\to\OO_W\), represent the lifts on a common
neighborhood, and shrink the polydisc.  This gives a closed embedding
\(i:U\hookrightarrow D\subset\C^d\) and sections
\(b_1,\ldots,b_d\in\mathscr B_k(U)\) reducing to the coordinate
restrictions.  Put \(Z=S_{k,U}\), the open subspace of \(S_k\) with
underlying set \(\phi^{-1}(U)\).  By definition,
\(\mathscr B_k(U)=H^0(Z,\OO_Z)\).

For a possibly nonreduced complex space, a tuple of holomorphic
functions defines a holomorphic map to affine space, functorially on
structure sheaves \cite[Theorem~1.1, p.~10-02]{Grothendieck1961TechniquesIII}.
Apply this theorem to the \(b_j\).  Their point values lie in \(D\),
so the map factors as
\[
 g:Z\longrightarrow D.
\]
The same tuple theorem identifies its actual analytic restriction to
the reduction with \(i\phi\), because the reduced coordinate functions
are those of \(i\phi\).  In particular \(|g|=i|\phi|\).  The
restriction \(\phi^{-1}(U)\to U\) is proper and \(i\) is closed, so
\(g\) is proper; nilpotents do not change properness of the underlying
continuous map.

The holomorphic map \(g\) makes
\(\mathscr A=g_*\OO_Z\) a genuine \(\OO_D\)-algebra.  Its underlying
module is coherent by the proper direct-image theorem
\cite[Theorem~1.1]{Grothendieck1961TechniquesVIII}.  Directly from the
underlying maps, as sheaves of \(\C\)-algebras,
\[
 \mathscr A=i_*(\mathscr B_k|_U).
\]
Reduction gives an epimorphism
\(\mathscr A\to i_*\OO_U\).  It is \(\OO_D\)-linear because the
actual analytic restriction \(g|_S=i\phi\) identifies the reduced
action of every ambient holomorphic function with its restriction to
\(U\).  The target is the coherent quotient of \(\OO_D\) by the
ideal of \(U\).  Hence its kernel \(\mathscr N\) is a coherent
\(\OO_D\)-module.  It is also a nilpotent ideal.  Shrink \(D\) once
more and choose sections \(n_1,\ldots,n_r\) generating
\(\mathscr N\) as a module.  They are actual sections of \(\OO_Z\)
whose reductions vanish.  Since \(\OO_D\to i_*\OO_U\) is an
epimorphism, subtraction of an ambient lift of a reduction gives the
stalkwise module equality
\begin{equation}\label{eq:analytic-module-generation}
 \mathscr A=\OO_D\cdot1+
              \sum_{\nu=1}^r\OO_D\cdot n_\nu.
\end{equation}
No splitting of the reduction map is used here.

Apply the tuple theorem again to the \(b_j\) and \(n_\nu\).  It gives
\[
 H=(g,n_1,\ldots,n_r):Z\longrightarrow P=D\times\C^r.
\]
If \(j:U\hookrightarrow P\) is \(u\mapsto(i(u),0)\), then
\(|H|=j|\phi|\), since the \(n_\nu\) are nilpotent.  Thus \(H\)
is proper.  The sheaf \(\mathscr A'=H_*\OO_Z\) is a coherent
\(\OO_P\)-module and the unit is an algebra map
\[
 \theta:\OO_P\longrightarrow\mathscr A'.
\]
As sheaves of rings \(\mathscr A'=j_*(\mathscr B_k|_U)\), so its
stalk at \(j(u)\) is \(\mathscr B_{k,u}=\mathscr A_{i(u)}\), and
its stalk outside \(j(U)\) is zero.  The stalk identification uses
the cofinality of restrictions of ambient neighborhoods among
neighborhoods in the closed subspace \(U\).

The map \(\theta\) is surjective on every stalk.  At \(j(u)\), any
element of the target has, by \eqref{eq:analytic-module-generation},
the form \(a_0+\sum a_\nu n_\nu\) with
\(a_\nu\in\OO_{D,i(u)}\).  The projection of \(H\) to \(D\) is
\(g\), and the extra coordinate \(t_\nu\) pulls back to \(n_\nu\).
The element is therefore the image of the convergent germ
\[
 a_0(z)+\sum_{\nu=1}^r a_\nu(z)t_\nu\in\OO_{P,j(u)}.
\]
At all other stalks the target is zero.  The kernel
\(\mathscr Q=\ker\theta\) is a coherent analytic ideal
\cite[pp.~9-07--9-08]{Grothendieck1961TechniquesII}.  It defines a closed
complex subspace \(Y=V(\mathscr Q)\subset P\) by
\cite[Definition~2.2 and adjacent construction, p.~9-10]{Grothendieck1961TechniquesII}.
Its support is exactly
\(j(U)\), since the target stalk there surjects onto the nonzero ring
\(\OO_{U,u}\).  Under the homeomorphism with \(U\), its structure
sheaf is \(\mathscr B_k|_U\), and its reduction is \(U\).

Every germ of \(\mathscr Q\) pulls back to zero by the definition of
\(\theta\).  The closed-subspace factorization criterion
\cite[pp.~9-04--9-05]{Grothendieck1961TechniquesII} therefore factors
\(H\) holomorphically through \(Y\).  Its sheaf map, under the
identified structure sheaf, is the canonical evaluation
\[
 |\phi|^{-1}(\mathscr B_k|_U)\longrightarrow\OO_Z.
\]
Indeed every germ of \(\mathscr B_k|_U\) lifts through the surjective
\(\theta\) to an ambient germ, whose pullback is precisely its value
as a section on the inverse image.  This proves the assertion for every
germ, not only the chosen coordinates.

\par\medskip\noindent\textit{Gluing and closed transition maps.}\par
The ringed space \((|W|,\mathscr B_k)\) was fixed before any of the
local choices.  The presentations just constructed make each of its
restrictions an analytic space by the local definition and adjacent
construction in \cite[Definitions~2.1--2.2 and adjacent prose]{Grothendieck1961TechniquesII}.
On overlaps the identifications are
the identity on the same sheaf of \(\C\)-algebras, and hence are
analytic isomorphisms satisfying the cocycle condition
\cite[Definition~2.4 and Proposition~2.5]{Grothendieck1961TechniquesII}.
This proves that \(W_k\) is a complex space.  It retains the Hausdorff
topology of \(W\), and its reduction is canonically \(W\).
Write \(\varepsilon_k:|\phi|^{-1}\mathscr B_k\to\OO_{S_k}\) for
the global sheaf evaluation.  It is locally the holomorphic factorization
above, so it defines \(\phi_k\).  Its reduction is \(\phi\), and its
underlying map is \(|\phi|\); thus it is proper and surjective.
The direct-image equality is the definition of its target structure
sheaf, with the canonical unit as the identification.

For the closed transition, use the construction with \(H\) at level
\(k+1\) and restrict its coordinate functions to the closed subspace
\(S_{k,U}\).  They give another proper holomorphic map to the same
\(P\), with underlying map \(j|\phi|\).  Its unit is the composite
\[
 \OO_P\twoheadrightarrow j_*(\mathscr B_{k+1}|_U)
          \twoheadrightarrow j_*(\mathscr B_k|_U).
\]
The second quotient stays surjective under the closed embedding \(j\),
as seen on its stalks.  The last sheaf with this ambient action is
coherent by proper direct image.  The two units therefore have nested
coherent ideal kernels, so their analytic subspaces give the required
closed inclusion.  These local inclusions glue because their ring maps
are the fixed quotients.  Naturality of the evaluation maps gives
\[
 \varepsilon_k\circ|\phi|^{-1}(\mathscr B_{k+1}\to\mathscr B_k)
  =(\OO_{S_{k+1}}\to\OO_{S_k})\circ\varepsilon_{k+1},
\]
which proves commutativity with the source inclusions.  No retraction
\(W_k\to W\), Cartesian square, or realization of the infinite tower
as one formal completion is asserted or needed.
\end{proof}

\subsection{Extension of a prime-floor contraction}
\label{subsec:prime-floor-extension}

\begin{proposition}[Extension from a prime floor]\label{prop:prime-floor-extension}
Under the hypotheses of Proposition~\ref{prop:conormal-layers}, there
are a normal compact analytic space \(Z\), a closed embedding
\(W\hookrightarrow Z\), and a proper bimeromorphic morphism
\(F:X\to Z\) such that \(F|_S\) is the given map \(\phi\) followed
by that embedding and
\begin{equation}\label{eq:extension-fibers}
 X\setminus S\simeq Z\setminus W,\qquad
 F^{-1}(|W|)=|S|,\qquad
 |F^{-1}(w)|=|\phi^{-1}(w)|\quad(w\in W).
\end{equation}
The natural map \(\OO_Z\to F_*\OO_X\) is an isomorphism, and the
fibers are connected.  For a sufficiently divisible global index
\(\ell\) of \(S\), the actual line \(\OO_X(-\ell S)\) is
\(F\)-ample; in particular \(F\) is projective.  More generally an
actual rational line on \(X\) whose restriction to \(S\) is
\(\phi\)-ample is \(F\)-ample.  The target is in Fujiki's class
\(\mathcal C\).
\end{proposition}
\begin{proof}
Choose a positive global index \(\ell\) such that \(\ell S\) is
Cartier and \(\OO_S(-\ell S)\) is \(\phi\)-ample.  Put
\[
 A=S_\ell=\ell S,\qquad A'=W_\ell,
 \qquad f=\phi_\ell:A\longrightarrow A'.
\]
Here \(A\) is the full effective Cartier divisor, including its
nilpotents; its local equation on normal \(X\) is a nonzerodivisor.
Lemma~\ref{lem:analytic-thickenings} makes \(f\) a proper surjective
holomorphic map of complex spaces and gives the natural equality
\(f_*\OO_A=\OO_{A'}\).

Cartier multiplication gives
\(\mathscr I_\ell\mathscr I_j=\mathscr I_{j+\ell}\).  Hence, for
every integer \(n>0\), on the full space \(A\) one has
\begin{equation}\label{eq:cartier-negative-layer}
 \OO_A(-nA)
 =\mathscr I_{n\ell}\otimes\OO_A
 =\mathscr I_{n\ell}/\mathscr I_{(n+1)\ell}.
\end{equation}
Its finite filtration by
\(\mathscr I_j/\mathscr I_{(n+1)\ell}\), for
\(n\ell\leq j\leq(n+1)\ell\), has successive quotients
\(\mathscr E_j\) for \(n\ell\leq j<(n+1)\ell\).  All indices are
at least one.  These filtration terms are \(\OO_A\)-modules, because
\(\mathscr I_\ell\mathscr I_j\subset\mathscr I_{(n+1)\ell}\).

Let \(a:S\hookrightarrow A\) and \(i:W\hookrightarrow A'\) be the
closed reduction embeddings.  The compatibility in
Lemma~\ref{lem:analytic-thickenings} gives \(fa=i\phi\).
Pushforward by a closed embedding is exact, so the composition identity
for derived direct images of abelian sheaves gives
\[
 R^q f_*(a_*\mathscr E_j)=i_*R^q\phi_*\mathscr E_j=0\quad(q>0).
\]
The long exact sequences of the finite filtration in
\eqref{eq:cartier-negative-layer} therefore give
\begin{equation}\label{eq:cartier-negative-vanishing}
 R^q f_*\OO_A(-nA)=0\qquad(q>0,\ n>0).
\end{equation}
This uses no vanishing for \(\mathscr E_0\) and keeps the full Cartier
thickening throughout.

The actual line \(\OO_A(-A)\) restricts to
\(\OO_S(-\ell S)\).  The underlying map of \(f\) is \(\phi\),
so the reductions of their fibers, as closed subspaces of \(X\),
coincide: a reduced closed analytic subspace is determined by its support.
Thus \(\OO_A(-A)\) is ample on the reduction of every \(f\)-fiber.
For a line on a compact nonreduced complex space, ampleness on the
reduction implies ampleness on the whole space.  Indeed the positive
metric from \cite[Corollary~1.12]{Fujino2026Fujiki} has local strictly
plurisubharmonic weights in common ambient embeddings, by its Lemma~2.4.
The same weights define a metric on the original line, since the
pointwise absolute values of transition units are unchanged by
nilpotents.  The converse positive-metric criterion gives ampleness on
the full space.  The proper fiberwise criterion
\cite[Definition~3.1 and Remark~3.2]{Fujino2026Fujiki} now makes
\(\OO_A(-A)\) \(f\)-ample.

We apply the analytic blowdown criterion stated in
\cite[Theorem~4.2, p.~14]{DasHacon2024OnMMP}, attributed there to
\cite[Theorem~2, p.~495]{Fujiki1975}.  For a reduced complex space,
an effective Cartier divisor \(A\) with its full possibly nonreduced
structure, and a proper surjective holomorphic map \(f:A\to A'\),
this criterion assumes \(f\)-ampleness of \(\OO_A(-A)\) and
\[
 R^1f_*\OO_A(-nA)=0\qquad\hbox{for every }n>0.
\]
It supplies a blowing down: a complex target \(Z\) containing \(A'\)
as an embedded subspace, a proper surjective holomorphic map
\(F:X\to Z\) whose restriction to \(A\) is \(f\), and an analytic
isomorphism of the complements.  In its universal form it also supplies
the natural equality of sheaves of rings
\begin{equation}\label{eq:universal-blowdown-ring}
 F_*\mathscr S_{X,A,f}=\OO_Z.
\end{equation}
Here \(\mathscr S_{X,A,f}\) consists along \(A\) of the germs in
\(\OO_X\) whose restrictions lie locally in
\(\operatorname{im}(f^{-1}\OO_{A'}\to\OO_A)\), and is \(\OO_X\)
away from \(A\).  This is a subsheaf of rings; no module coherence of
this particular sheaf is used.  The hypotheses of the stated criterion
follow from \eqref{eq:cartier-negative-vanishing} and the preceding
ampleness argument.  No flatness or reducedness of \(A\) or \(A'\)
is required.

The embedded \(A'\) is closed because it is compact.  Its reduction
embeds the original \(W\) as a closed subspace of \(Z\).  The
compatible reductions give \(F|_S=\phi\) followed by this embedding.
Removing a subspace depends only on its support, so the complement
isomorphism is \(X\setminus S\simeq Z\setminus W\).  It follows
that \(F^{-1}(|W|)=|S|\), and the prescribed restriction to \(A\)
gives the fiber-set equalities in \eqref{eq:extension-fibers}.
Outside \(W\) the fibers are single reduced points.  These assertions
concern underlying fiber sets over \(W\); they require no Cartesian
square of the closed thickenings.  The fibers are connected because
the floor fibers are connected.  The target is compact as the image of
compact \(X\).

We next prove normality; it is not an extra conclusion assumed from the
blowdown criterion.  For every open \(V\subset Z\), a section
\(s\in\OO_X(F^{-1}V)\) restricts to a section on the open space
\[
 A\cap F^{-1}V=f^{-1}(V\cap A').
\]
The equality \(f_*\OO_A=\OO_{A'}\) makes this restriction the pullback
of a unique section on \(V\cap A'\).  Thus every germ of \(s\) along
\(A\) satisfies the defining condition of \(\mathscr S_{X,A,f}\).
The reverse inclusion follows from \(\mathscr S_{X,A,f}\subset\OO_X\).
Consequently
\[
 F_*\mathscr S_{X,A,f}=F_*\OO_X,
 \qquad F_*\OO_X=\OO_Z
\]
as sheaves of rings, the second equality by
\eqref{eq:universal-blowdown-ring}.  This reasoning uses equality of
the displayed open subspaces, not of any closed fiber products.

The equality first makes \(Z\) reduced.  The complement is dense:
the inverse image of a nonempty open of \(Z\) is a nonempty open of
\(X\), and therefore meets \(X\setminus S\).  The complement
isomorphism thus makes \(F\) bimeromorphic.  Let
\(\nu:Z^\nu\to Z\) be the finite analytic normalization.  The
finite normalization exists by
\cite[Part~B, Section~4, Corollaries~2--3]{Houzel1961}.  The
dominant bimeromorphic map from normal \(X\) lifts to \(Z^\nu\).
Pullback along this lift gives inclusions, injective on the dense common
complement,
\[
 \OO_Z\subset\nu_*\OO_{Z^\nu}\subset F_*\OO_X=\OO_Z.
\]
The normalization is consequently an isomorphism and \(Z\) is normal.
The natural direct-image equality also gives connected fibers by analytic
Stein factorization.

Finally consider the actual global line \(\OO_X(-\ell S)\).  Its
restriction is ample on every reduced \(F\)-fiber: this is clear off
\(W\), and over \(W\) the reduction is the same embedded reduced
space as that of the corresponding \(\phi\)-fiber.  The nonreduced
upgrade and the proper fiberwise criterion used above therefore make
this one line \(F\)-ample.  The same proof works for any actual ambient
rational line with \(\phi\)-ample restriction after clearing its
global index.  In particular it works for the negative full adjoint
in Proposition~\ref{prop:conormal-layers}.  Thus \(F\) is projective.
A projective resolution with smooth source of the compact K\"ahler
\(X\) is again compact K\"ahler, and its composite with \(F\) is
bimeromorphic.  Hence \(Z\) lies in Fujiki's class \(\mathcal C\).
\end{proof}

The proposition supplies the ambient step for floor data satisfying its
stated hypotheses, including the two actual ample-line conditions.  In
the threefold argument below it extends divisor-to-curve contractions.  In
Section~\ref{sec:reduction} it extends a selected floor contraction of
the supported fourfold program.  When that floor contraction cuts out
the selected ray, the exact fiber description keeps the same contracted
curves, and the actual relatively ample lines make the ambient step
projective.  Its target K\"ahler class is then established in the
supported-program argument.


\subsection{Descent of a dominating current}

\begin{lemma}[Descent of a dominating current]
\label{lem:dominating-current-descent}
Let \(p:Y\to Z\) be a proper bimeromorphic morphism between normal
compact complex spaces, with \(Y\) K\"ahler.  Let \(\eta\) be a smooth
real closed \((1,1)\)-form with local smooth potentials on \(Z\).
Suppose that a positive closed current \(T\) with local potentials
represents the local-potential Bott--Chern class \(p^*[\eta]\), and that
\(T\geq\kappa\) for a K\"ahler form \(\kappa\) on \(Y\).  For every
smooth positive Hermitian form \(g\) on \(Z\), understood in local
ambient embeddings, there are a constant \(c>0\) and a global
quasi-plurisubharmonic function \(u_Z\) such that
\begin{equation}\label{eq:dominating-current-descent}
 \eta+i\partial\bar\partial u_Z\geq c g.
\end{equation}
This is a current in \([\eta]\) with local potentials, and it has the
bigness property preceding \cite[Theorem~2.29]{DasHaconPaun2024}.
\end{lemma}
\begin{proof}
We make the potential in the class equality explicit.  Use the
normal-space potential description of Boucksom--Guedj
\cite[Section~4.6.1]{BoucksomGuedj2013}, specifically Lemma~4.6.1 and
the paragraph following Definition~4.6.2.  The soft-sheaf descriptions
of Bott--Chern cohomology by smooth functions and distributions first
give a global real distribution \(u_Y\) for which
\[
 T=p^*\eta+i\partial\bar\partial u_Y.
\]
Locally write \(T=i\partial\bar\partial v\) with \(v\) plurisubharmonic
and \(p^*\eta=i\partial\bar\partial h\) with \(h\) smooth.  Then
\(u_Y+h-v\) is pluriharmonic as a distribution.  By the cited lemma
it is a smooth real part of a holomorphic germ.  Thus \(u_Y\) is locally
represented by a quasi-plurisubharmonic function.  These representatives
agree as distributions on overlaps and hence, after taking their
upper-semicontinuous representatives, agree pointwise.  They give one
global quasi-plurisubharmonic \(u_Y\), which is locally integrable and
locally bounded above.

In local ambient embeddings, \(p^*g\) is represented by a smooth
semipositive form and \(\kappa\) by a smooth positive definite form.
A finite relatively compact cover of \(Y\) therefore gives a constant
\(C>0\) with \(p^*g\leq C\kappa\).  Consequently
\[
 p^*\eta+i\partial\bar\partial u_Y\geq C^{-1}p^*g.
\]
The map \(p\) is a modification between normal spaces, and \(u_Y\)
has the local upper bound needed in the current descent argument.
\cite[Corollary~2.32 and its proof, p.~18]{DasHaconPaun2024} gives a global
quasi-plurisubharmonic \(u_Z\) with
\eqref{eq:dominating-current-descent} for \(c=C^{-1}\).  It is an
\(L^1\) potential, and the formula keeps the descended current in the
chosen Bott--Chern class.  The local potentials of \(\eta\) give local
potentials for that current.  No K\"ahler form on \(Z\) has been used.
\end{proof}

\subsection{Threefold contractions used in the boundary argument}
\label{subsec:threefold-contraction-bridge}

We prove Proposition~\ref{prop:threefold-contraction} by applying the
cited threefold MMP with two analytic constructions supplied here:
Proposition~\ref{prop:finite-null} contracts the finite null loci, and
Proposition~\ref{prop:prime-floor-extension} extends the prime-floor
contractions.  The projective relative MMP, the nonbig argument through
a surface contraction, and the terminal canonical threefold program
remain external inputs.  We verify the hypotheses at each use of the
two constructions, then supply the required target positivity; neither
construction alone asserts that its target is K\"ahler.  We first record
the cone observation used in the big and floor cases.

For a normal compact space \(Z\) in Fujiki's class \(\mathcal C\), put
\(N^1(Z)=H^{1,1}_{BC}(Z)\) in the local-potential convention.  Let
\(N_1(Z)\) be the vector space of real closed currents of bidimension
\((1,1)\), modulo \(T_1\equiv T_2\) when \(T_1(\eta)=T_2(\eta)\) for every
real closed \((1,1)\)-form \(\eta\) with local potentials; the pairing is
evaluation.  Following
\cite[Definitions~3.6 and~3.8, pp.~8--9]{HoeringPeternell2016}, let
\(\operatorname{NA}(Z)=\overline{\operatorname{NA}}(Z)\subset N_1(Z)\)
denote the closed cone generated by the numerical classes of positive
closed currents.  Thus both notations used below refer to this same closed
cone.

\begin{lemma}[Positivity on the current cone]\label{lem:na-slice}
Let \(Z\) be a normal compact K\"ahler analytic variety with rational
singularities, and let \(\kappa\) be a K\"ahler class.  In the
finite-dimensional local-potential numerical spaces put
\[
 C=\overline{\operatorname{NA}}(Z),\qquad
 \Sigma=\{z\in C:\kappa\cdot z=1\}.
\]
Then \(\Sigma\) is compact.  If a real local-potential \((1,1)\)-class
\(\ell\) is strictly positive on \(\Sigma\), then \(\ell\) is
K\"ahler.
\end{lemma}
\begin{proof}
By \cite[Lemma~2.3 and Proposition~2.4]{DasHaconPaun2024}, the natural
pairing is perfect, \(\operatorname{Nef}(Z)=C^*\), and the K\"ahler
cone is open with closure the nef cone.  In particular \(\kappa\) lies
in the interior of \(C^*\).  A sequence in \(\Sigma\) with unbounded
norm would, after division by its norm and passage to a subsequence,
give a nonzero class \(z\in C\) with \(\kappa\cdot z=0\), contradicting
that interior property.  The slice is closed and hence compact.

If it is nonempty, let \(\mu=\min_\Sigma\ell>0\).  Homogeneity gives
\(\ell-\frac\mu2\kappa\in C^*=\operatorname{Nef}(Z)\).  A smooth
local-potential nef approximation for this last class, bounded below by
\(-\mu\kappa/4\), becomes a positive representative after adding
\(\mu\kappa/2\).  Thus \(\ell\) is K\"ahler.  If \(\Sigma\) is
empty, then \(C=\{0\}\) by the interior property; the same conclusion
follows because every class is nef and a positive small multiple of
\(\kappa\) can be subtracted first.  All cones here are the
local-potential cones in the cited perfect pairing.
\end{proof}

The spaces used below have rational singularities: this is Lemma~2.31
of \cite{DasHacon2024} for dlt pairs and its Remark~2.14 for klt varieties.

\begin{proof}[Proof of Proposition~\ref{prop:threefold-contraction}]
The proof of
\cite[Theorem~5.5]{DasHaconPaun2024} first uses a projective relative
small \(\Q\)-factorialization and then
\cite[Theorem~1.7]{DasHacon2024}.  The latter proof passes to a small
strongly \(\Q\)-factorial model and treats a nef nonbig class by its
Theorem~5.5 and a nef big class by its Theorem~6.4.  Here strong
\(\Q\)-factoriality means that every global coherent reflexive rank-one
sheaf has an invertible reflexive power.  It implies ordinary global
\(\Q\)-factoriality; it does not assert the analogous property for
all analytic open subsets.  The small models use the projective relative
MMP of \cite[Proposition~2.26 and Lemma~2.27]{DasHacon2024}, whose maps
are already projective.  Compactness of the final bases is retained when
projective maps are composed, as required by
\cite[Remark~2.11]{DasHaconPaun2024}.

\paragraph{The imported programs and target descent.}
The nef nonbig branch follows the chain Theorem~5.5, Corollary~5.4,
and Theorem~5.2 of \cite{DasHacon2024}.  It uses Step~4 in the proof of
\cite[Theorem~1.4, p.~242]{HoeringPeternell2015}, which contracts a
negative-definite collection of curves on a smooth compact K\"ahler
surface and descends through a graph.  We use that nonbig argument as
an external input.  The pseudo-effective cone lineage also uses the
terminal canonical threefold program: Assumption~10.1 of
\cite{DasOu2022} invokes its nonvanishing Theorem~9.1, whose proof at
p.~50 invokes \cite[Theorem~1.1]{HoeringPeternell2015} for a terminal
\(K\)-MMP and Mori fiber space.  We retain that terminal program as
an external input.  Together with the projective relative MMP just
specified, these are the birational inputs used below.

For an already constructed negative-ray contraction with ordinary
\(\Q\)-factorial compact K\"ahler dlt source and relatively ample
negative adjoint, \cite[Proposition~3.1]{CampanaHoeringPeternell2016}
gives rationality of the target and descent of the supporting class.
We use the proof of its Corollary~3.1 for target positivity.  The
corollary assumes non-uniruledness to obtain a nef support; our ray
already has an exposed nef support.  Its remaining argument descends
that class, proves positivity on the target current cone, and lifts
target curves through the projective morphism.  These steps apply to
the data just stated.  The proposition and corollary numbers here are
those of the journal edition.  We now verify the finite-null and
prime-floor constructions to which this argument will be applied.

\paragraph{Nef and big classes with a finite null locus.}
Three geometric situations use Proposition~\ref{prop:finite-null}:
a small negative ray, the entire null locus at a stopping model, and a
small ray with pseudo-effective adjoint.  The big threefold program uses
strongly \(\Q\)-factorial klt sources.  The last situation can also
occur for ordinary dlt data, whose underlying space is ordinary
\(\Q\)-factorial and klt, as checked below.  In every case the relevant
locus is the full top-intersection null locus of the proposition.

For a small negative ray, the small branch of
\cite[proof of Theorem~4.16, p.~40]{DasHacon2024} has, after its boundary
perturbation, a strongly \(\Q\)-factorial klt source and a nef big
supporting class \(a=(K+\Gamma')+\kappa\) with \(\kappa\) K\"ahler.
At this input \cite[Lemma~4.5]{DasHacon2024} says that an \(a\)-null
surface is covered by \(a\)-null curves.  All those curves lie in the
exposed ray.  Such a covering of a surface contradicts the smallness of
that ray.  Lemma~\ref{lem:no-null-surface} therefore makes
\(\operatorname{Null}(a)\) a finite union of curves, and
Proposition~\ref{prop:finite-null} constructs its contraction.
This argument permits a non-pseudo-effective underlying adjoint.

For the entire null locus at the final model \(X_n\) of the program in
\cite[Claim~6.5 and proof of Theorem~6.4, pp.~50--51]{DasHacon2024},
the pair is strongly \(\Q\)-factorial klt, its class \(a_n\) is nef
and big, and every \(a_n\)-null curve has nonnegative degree for the
running adjoint \(A_n\).  These are the stopping data from that program.
Its null-surface argument, using Lemma~4.5 of the same source, would
cover a null surface by \(a_n\)-null curves on which the remainder
\(a_n-c_1(A_n)\) has positive degree.  Their \(A_n\)-degrees would
then be negative, contrary to the stopping data.  There is therefore no
null surface.  Lemma~\ref{lem:no-null-surface} again gives finitely many
null curves.  Apply Proposition~\ref{prop:finite-null} to their entire
union.  This use has no single-ray hypothesis and needs only the proper
bimeromorphic map to a normal compact analytic target supplied by that
proposition.

For a small ray with pseudo-effective adjoint, the branch in the proof of Theorem~6.4 at
p.~50 invokes the small case of Theorem~2.23(1) of the same source;
this case can also occur in its Theorem~3.1, Claim~3.2.  Let \(A\) be
the negative adjoint and \(b\) an exposed nef support for the ray.
It may be scaled so that \(b-c_1(A)\) is K\"ahler.  To check this
usual support step, normalize the closed cone of positive
bidimension-\((1,1)\) classes by a K\"ahler class.  Its normalized slice
is compact.  The continuous function \(-c_1(A)\) is strictly positive
on its intersection with the exposed ray, hence on a neighborhood of
that intersection.  On the compact complement, the support has a
positive minimum.  A large multiple of the support minus \(c_1(A)\)
is consequently strictly positive on the full slice.
Lemma~\ref{lem:na-slice} makes this difference K\"ahler.  Since
\(A\) is pseudo-effective, the scaled
\(b\) is big.  The same Lemma~4.5 and the smallness of the ray exclude
null surfaces, so Lemma~\ref{lem:no-null-surface} and
Proposition~\ref{prop:finite-null} apply.  At the outer Theorem~6.4
use the pair is klt.  If this step is made for the ordinary dlt data in
Claim~3.2, the underlying space is still klt, which is the singularity
condition of Proposition~\ref{prop:finite-null}; a small decrease of
the boundary preserves negativity when a klt adjoint is needed in a
subsequent result.

In each of the two single-ray situations just described, the constructed
map contracts exactly that ray.  A positive-dimensional reduced fiber
is a finite connected union of compact curves, hence is projective.
For the corresponding negative adjoint \(A\), a positive global Cartier
multiple of the actual line \(-A\) has positive degree on every
irreducible component of that fiber, and is
therefore ample on the fiber.  It is ample also on every zero-dimensional
fiber.  Ampleness passes from a reduction to its nilpotent thickening by
the positive-metric criterion
\cite[Lemma~2.4 and Corollary~1.12]{Fujino2026Fujiki}.  The proper
fiberwise criterion \cite[Definition~3.1 and Remark~3.2]{Fujino2026Fujiki}
now makes this one actual line relatively ample, so the contraction is
projective.  Normality and proper bimeromorphicity give connected fibers.
The source is ordinary \(\Q\)-factorial klt, or dlt in the indicated
variant, and the exposed nef support is unchanged.  Thus
\cite[Proposition~3.1 and proof of Corollary~3.1]{CampanaHoeringPeternell2016},
applied as specified above, gives the K\"ahler target and the descended
supporting class.
Thus Proposition~\ref{prop:finite-null} supplies the normal compact
analytic contraction, and this postprocessing supplies projectivity and
positivity in the single-ray cases.  The entire-null cleanup uses the
different argument that follows.

In that cleanup, write \(\mu:X_n\to Z\) for the entire-null contraction
just constructed.  The graph argument in
\cite[proof of Theorem~6.4, pp.~51--52]{DasHacon2024} descends its
composite with the program to a morphism \(\psi:X_0\to Z\), where
\(X_0\) is the original strongly \(\Q\)-factorial klt source of that
theorem.  All steps of the program are trivial for the transported
supporting class.  Before applying its Lemma~2.44 to \(\psi\), restore
the original boundary \(\Gamma^{\rm orig}\) and the original nef big
difference \(\omega^{\rm orig}\) on \(X_0\), before the internal
boundary replacement.  The original supporting class is
numerically trivial on its contracted curves, and hence
\begin{equation}\label{eq:cleanup-original-boundary}
 -(K+\Gamma^{\rm orig})\equiv_\psi\omega^{\rm orig}.
\end{equation}
The left side is \(\psi\)-nef and is \(\psi\)-big for the bimeromorphic
map.  The rationality result of Lemma~2.44 thus applies to this original
Q-Gorenstein klt pair and makes the target \(Z\) rational.  This step
uses the original nef big difference; the transported remainder
\(a_n-c_1(A_n)\) on the stopping model is not substituted into
\eqref{eq:cleanup-original-boundary}.

For the entire-null contraction \(\mu\), the source \(X_n\) is a
normal compact K\"ahler klt
threefold, so it has rational singularities.  The target is normal and
compact, and is in Fujiki's class \(\mathcal C\): a projective resolution
of \(X_n\) with smooth source is compact K\"ahler, and its composite
with \(\mu\) is bimeromorphic.  The exceptional image
\(E_Z=\mu(\operatorname{Null}(a_n))\) is finite.  Every curve contracted
by \(\mu\) is \(a_n\)-null.  Apply
\cite[Lemma~2.11]{DasHacon2024} to this map between normal compact
rational spaces in class \(\mathcal C\).  It gives a smooth real closed
\((1,1)\)-form \(\eta_Z\) with local smooth potentials such that
\begin{equation}\label{eq:cleanup-descended-class}
 a_n=\mu^*[\eta_Z]
 \quad\text{in the local-potential Bott--Chern group.}
\end{equation}
The normalized graph constructed \(\psi\) earlier; this application of
Lemma~2.11 uses \(\mu\).

We check the three analytic positivity hypotheses on \(Z\).  Fix a
smooth positive Hermitian form \(g_Z\) in local ambient embeddings.
The bimeromorphic nef descent theorem
\cite[Lemma~3.13]{HoeringPeternell2016} applies to the normal compact
threefold \(Z\) in class \(\mathcal C\) and the normal compact
threefold \(X_n\).  It makes \([\eta_Z]\) analytically nef from
\eqref{eq:cleanup-descended-class}; it does not assume that \(Z\) is
K\"ahler.  If its defining approximation uses a different positive
reference form, compact comparison with \(g_Z\) and rescaling the
approximation parameter give, for every \(\delta>0\), a smooth
function \(f_\delta\) with
\begin{equation}\label{eq:cleanup-nef}
 \eta_Z+i\partial\bar\partial f_\delta\geq-\delta g_Z.
\end{equation}
Changing a smooth representative of the same local-potential class only
changes \(f_\delta\) by a global smooth potential.

The big class \(a_n\) on the compact K\"ahler source contains a positive
closed current \(T_n\) with local potentials and
\(T_n\geq\kappa_n\) for a K\"ahler form \(\kappa_n\).  In view of
\eqref{eq:cleanup-descended-class}, Lemma~\ref{lem:dominating-current-descent}
applied to \(\mu\) gives a global quasi-plurisubharmonic, hence
\(L^1\), function \(u_Z\) and a constant \(c>0\) satisfying
\begin{equation}\label{eq:cleanup-big}
 \eta_Z+i\partial\bar\partial u_Z\geq c g_Z.
\end{equation}
This establishes the required current bigness on the not yet K\"ahler
target in the selected class \([\eta_Z]\).

Finally let \(V\subset Z\) be an irreducible reduced analytic subspace
of dimension \(d>0\).  Its strict transform \(V'\) is the closure
of the inverse image of \(V\setminus E_Z\).  It has dimension \(d\),
maps bimeromorphically to \(V\), and is not contained in
\(\operatorname{Null}(a_n)\).  Nef approximations on \(X_n\), their
restriction to the integration cycle of \(V'\), and Stokes' theorem
give \(a_n^d\cdot V'\geq0\).  Equality would put \(V'\) into the
null locus by its definition, so the inequality is strict.  The proper
cycle identity \(\mu_*[V']=[V]\) and projection formula therefore give
\begin{equation}\label{eq:cleanup-positive-intersections}
 \int_{V_{\rm reg}}\eta_Z^d=a_n^d\cdot V'>0.
\end{equation}
This includes \(V=Z\); it requires neither \(V\) nor \(V'\) to be
normal.  Additivity over the top-dimensional irreducible components
gives the same positivity for every positive-dimensional compact reduced
analytic subspace.

Equations \eqref{eq:cleanup-nef}, \eqref{eq:cleanup-big}, and
\eqref{eq:cleanup-positive-intersections} are the three hypotheses of
\cite[Theorem~2.29]{DasHaconPaun2024}.  Apply that theorem directly to
\(\eta_Z\).  It gives a smooth potential making \(\eta_Z\) positive
definite, so \([\eta_Z]\) is K\"ahler and \(Z\) is compact K\"ahler.
This verifies the positivity conclusion at the cleanup contraction.

\paragraph{Divisorial contractions.}
In the point case, let \(P\) be the selected prime surface and
\(\nu:P^\nu\to P\) its normalization.  The nef-dimension-zero
condition for the supporting class is taken on \(P^\nu\); the whole
normalized fiber has zero restricted class by
\cite[Theorem~3.19(a) and its proof, p.~234]{HoeringPeternell2015}.
Use \cite[Corollary~4.4]{DasHacon2024OnMMP}, whose proof invokes its
Lemma~4.3, at the stated data: a \(\Q\)-factorial compact K\"ahler
source, a nef big exposed support, ray curves covering \(P\), and zero
restriction on \(P^\nu\).  The corollary constructs the proper analytic
point contraction and its proof supplies an actual ample line
\(\OO_P(-mP)\) for some \(m>0\).  Thus \(P\) is projective, and
ampleness persists on the full possibly nonreduced fiber by the
positive-metric criterion used above.  The full-fiber criterion makes
\(-mP\) relatively ample.

Every curve of \(P\) lifts finitely to its normalization.  The zero
restricted support therefore gives zero degree on every such curve, so
its ambient class lies in the exposed ray.  The negative adjoint
restricted to \(P\) is consequently numerically a fixed positive
multiple of the negative normal line.  Both have actual rational Cartier
multiples.  Numerical invariance of ampleness on the projective \(P\)
and then the full-fiber criterion make the negative adjoint relatively
ample.  The same
\cite[Proposition~3.1 and proof of Corollary~3.1]{CampanaHoeringPeternell2016}
then supplies rationality, class descent, and target positivity for this
negative-ray contraction.  The point-contraction corollary was used
to construct the underlying proper analytic map.

For a divisor-to-curve step, Proposition~\ref{prop:prime-floor-extension}
supplies the ordinary extension used both in
\cite[Theorem~3.1, Claim~3.2]{DasHacon2024} and in its replay in
Theorem~4.16 of that source.  We verify the two ample-line hypotheses
for this use.  Write \((M,\Delta_M)\) for the strongly
\(\Q\)-factorial compact K\"ahler dlt threefold there,
\(A=K_M+\Delta_M\), \(P\) for the selected prime floor, and \(R\)
for its exposed ray.  The selected ray satisfies
\[
 A\cdot R<0,\qquad P\cdot R<0.
\]
Adjunction makes \(P\) a normal compact K\"ahler dlt surface and makes
\(A|_P\) its actual full adjoint.  For its inclusion \(i:P\to M\), put
\[
 F=\{z\in\overline{\operatorname{NA}}(P):i_*z\in R\}.
\]
Fix a K\"ahler class \(\kappa\) on \(M\) and a nonzero \(r\in R\).
For \(z\in F\) with \(\kappa|_P\cdot z=1\), positivity of
\(\kappa|_P\) gives \(i_*z=r/(\kappa\cdot r)\).  Thus on the full
compact normalized vertical face
\begin{equation}\label{eq:inner-floor-positive-face}
 -(A|_P)\cdot z=\frac{-A\cdot r}{\kappa\cdot r}>0,
 \qquad
 -(P|_P)\cdot z=\frac{-P\cdot r}{\kappa\cdot r}>0 .
\end{equation}
This uses the cone of positive current classes and includes its limits.

Let \(\zeta\) be the nef support exposing \(R\).  Its restriction has
null cone \(F\).  The compact-slice argument just used makes
\(\lambda\zeta|_P-c_1(A|_P)\) K\"ahler for sufficiently large
\(\lambda\): near \(F\) the negative adjoint has the uniform positive
bound in \eqref{eq:inner-floor-positive-face}, and away from \(F\)
the support has a positive minimum.  Lemma~\ref{lem:na-slice} on the
rational compact K\"ahler surface turns this bound into a K\"ahler
class.  The ordinary dlt case of the surface theorem
\cite[Theorem~2.32]{DasHaconYanez2026} therefore gives a projective
surjection \(\gamma:P\to W_P\) with connected fibers onto a normal
compact K\"ahler space, contracting exactly \(F\), with the support
pulled back from a K\"ahler class on \(W_P\).  This use is of the
surface theorem.

Applying the same compact-slice estimate after adding a sufficiently
large pullback of that target K\"ahler class makes each of
\(-A|_P\) and \(-P|_P\) relatively K\"ahler over \(W_P\).
They have actual global rational Cartier multiples by the factoriality
of \(M\).  The proper fiberwise positivity criterion makes both lines
\(\gamma\)-ample, including when a fiber is the whole surface.
Put \(H=\floor{\Delta_M}-P\).  For small rational \(\varepsilon>0\),
the pair \((M,\Delta_M-\varepsilon H)\) is plt with sole floor \(P\).
Indeed on a dlt resolution the pullback of the effective globally
\(\Q\)-Cartier divisor \(H\) is effective.  Subtracting it preserves
the strict inequality for every exceptional crepant coefficient and
lowers all the other strict floor coefficients below one.  The negative
degree of \(A-\varepsilon H\) on \(R\) persists for small
\(\varepsilon\), so the first relative ample line persists by
\eqref{eq:inner-floor-positive-face}; the second is unchanged.
These are precisely the hypotheses of
Proposition~\ref{prop:prime-floor-extension}.

That proposition constructs the extension, whose fiber sets show that the ambient map
contracts exactly the original ray and is an isomorphism off \(P\).
The negative original adjoint is relatively ample by the same full-fiber
criterion.  Return at once to \(\Delta_M\) and \(A\) for the running
program.  For this projective negative-ray contraction,
\cite[Proposition~3.1 and proof of Corollary~3.1]{CampanaHoeringPeternell2016}
applies with the already exposed support, exactly as specified at the
start of the proof.  It gives rationality, class descent, and K\"ahler
target positivity.  Strong factoriality is preserved by
\cite[Lemma~2.5]{DasHacon2024}, and an
ordinary projective flip, when needed, is supplied by its Theorem~2.24.
Thus the original support and boundary are the ones transported to the
next step.

\paragraph{Return to the original threefold.}
The finite-null and divisorial constructions now supply the indicated
steps of the cited big-case program, with their required positivity.
Together with the external inputs stated at the start, they give the
asserted contraction for the data in
\eqref{eq:threefold-exposed-input} by the proof of
\cite[Theorem~1.7]{DasHacon2024}.  The proof of
\cite[Theorem~5.5]{DasHaconPaun2024} descends the contraction from the
small model.  A pulled-back K\"ahler class has degree zero on a constant
curve and positive degree on a nonconstant curve, as seen on its
normalization.  This proves the asserted curve criterion and completes
Proposition~\ref{prop:threefold-contraction}.
\end{proof}

For the separate exposure input, the proof of Theorem~5.2 and
Corollary~5.3 of \cite{DasHaconPaun2024} uses the same Theorem~1.7 for
the supporting contractions, followed by the projective relative cone
theorem and convex separation.  Its contraction inputs are therefore
the ones just established.

\subsection{The actual floor split}
\label{subsec:actual-floor-split}

We verify the use of that specialization for the dlt floor in the
fourfold supported program.  Let \((X,S+B)\) be one of its compact
ordinary \(\Q\)-factorial K\"ahler dlt pairs, with reduced floor
\(S\), coefficients of \(B\) below one, and actual adjoint
\(A=K_X+S+B\).  Fix a prime \(T\) of \(S\), and write
\[
 H=S-T,\qquad A_0=A-H=K_X+T+B.
\]
The finite global divisor \(H\) is \(\Q\)-Cartier by ordinary global
factoriality, and \(A_0\) is an actual rational line.  On a dlt
resolution, subtracting the effective pullback of \(H\) lowers the
exceptional crepant coefficients and removes the strict transforms of
the other floor components.  Only the strict transform of \(T\) has
coefficient one.  Hence \((X,T+B)\) is plt.  Ordinary plt adjunction
\cite[Lemma~5.3]{DasHacon2024OnMMP} gives a normal \(T\), an effective
rational boundary \(B_0\) with \((T,B_0)\) klt, and the actual
restriction \(A_0|_T=K_T+B_0\).

Choose one positive integer \(m\) clearing the indices of \(A,A_0,H\).
The canonical section of \(\OO_X(mH)\) restricts nontrivially to \(T\)
because \(T\) is not a component of \(H\).  Set
\[
 B'=\frac1m\operatorname{div}(s_{mH}|_T)\geq0.
\]
It is an effective rational \(\Q\)-Cartier divisor.  Use compatible
local meromorphic residue embeddings for the two adjunctions.  Their
ratio in codimension one on \(T\) is multiplication by \(s_{mH}|_T\).
Reflexive extension on normal \(T\) then gives the divisor split and
actual line identities
\begin{equation}\label{eq:actual-floor-boundary-split}
 \begin{aligned}
  B_T&=B_0+B',\\
  \OO_T(m(K_T+B_0))&\simeq\OO_X(mA_0)|_T,\\
  \OO_T(mB')&\simeq\OO_X(mH)|_T.
 \end{aligned}
\end{equation}
where \(A|_T=K_T+B_T\) is the full dlt adjunction.  This proves the
split without assuming \(T\) to be \(\Q\)-factorial.

Suppose, as in the supported program, that
\(\alpha_T=c_1(A|_T)+[\omega_T]\) is nef and \(\omega_T\) is
K\"ahler.  These properties hold by restricting the ambient smooth
metric nef approximations and local strictly plurisubharmonic potentials.
For a sufficiently small rational \(\varepsilon>0\), put
\[
 B_\varepsilon=B_0+(1-\varepsilon)B',\qquad
 \omega_\varepsilon=\omega_T+\varepsilon\theta_{B'},
\]
where \(\theta_{B'}\) is the normalized curvature of a smooth metric
on a Cartier multiple of \(B'\).  The full adjunction is lc and the
lower adjunction is klt.  Affineness of discrepancies in this convex
combination makes \((T,B_\varepsilon)\) klt.  On finitely many compact
local embedding charts, the Hessian of \(\theta_{B'}\) is bounded
and that of \(\omega_T\) is uniformly positive.  For the chosen small
\(\varepsilon\), \(\omega_\varepsilon\) is therefore K\"ahler and
\[
 \alpha_T=c_1(K_T+B_\varepsilon)+[\omega_\varepsilon].
\]
If \(\dim T=3\), Proposition~\ref{prop:threefold-contraction} applies
and gives a projective
connected-fiber contraction \(T\to W\) to a normal compact K\"ahler
target with \(\alpha_T\) pulled back from a K\"ahler class.  This is
the use of \cite[Corollary~5.6]{DasHaconPaun2024} in its Theorem~7.2.
The restricted class \(\alpha_T\) is not assumed big, and its null
face may have several rays.  The contraction follows from the assembled
argument in the preceding subsection, using
Proposition~\ref{prop:prime-floor-extension} and the specified external
results.
